\scanpage{11}

14) Supposons $\mathcal M$ semi-simple \inserted{sur $K\subset\mathbb R$} \unsure{engendrée|engendré} par motifs $M_1\ldots M_r$ de degrés $d_1\ldots d_r$, \inserted{et $K(1)$ \unsure{resp.|,} $K(-1)$} \deleted{\unsure{de}} \deleted{\unsure{objets}}. Pour tt $i$, soit $\varphi_i$ forme $(-1)^{d_i}$-sym. sur $M_i$, qui soit \deleted{\unsure{déf.} $\unsure{M_i}$} forme de Weil. \deleted{\unsure{et}} Pour qu'il existe une polarisation de $\mathcal M$ \deleted{définie} pour laquelle les $\varphi_i$ sont des formes de polarisation, il \deleted{\unsure{s}} faut qu'on ait la condition suivante :

a) Pour deux systèmes d'entiers $\underline m=(m_1,\ldots,m_{r+1}),\underline m'=(m'_1,\ldots,m'_{r+1})$ $\geq0$, avec $\overbrace{m_{r+1}+\sum_1^r m_i d_i}^{\delta(\underline m)}=\overbrace{m'_{r+1}+\sum m'_i d'_i}^{\delta(\underline m')}$, posant

\[
M_{\underline m}=\left(\bigotimes_{1\leq i\leq r}\left(\bigotimes^{m_i} M_i\right)\right)(\underset{\unsure{2|?}}{\unsure{+|-}}m_{r+1}),\quad
\varphi_{\underline m}=\left(\bigotimes_{1\leq i\leq r}\left(\underset{\deleted{\unsure{\varphi}}}{\bigotimes^{m_i}}\varphi_i\right)\right)(\underset{\unsure{2|?}}{\unsure{-|+}}m_{r+1}),\ \deleted{\unsure{\alpha}}
\]

on a, pour tout $u\in\operatorname{Hom}(M_{\underline m},M_{\underline m'})$, la relation

\[
\operatorname{Tr}u^{*(\varphi_{\underline m},\varphi_{\underline m'})}u\geq0
\]

Je dis que cette condition est aussi suffisante. \deleted{\unsure{Si elle est}} \inserted{Supposons celà en effet} \unsure{remplie|rempli}, soit $M\in\operatorname{Ob}\mathcal M^k$, \inserted{$M\ne0$.} par hyp. $\exists\ \underline m$ avec un $\alpha:M\hookrightarrow M_{\underline m}$ ($k=\deleted{\text{deg}}\sum m_i d_i$), soit $P(M)$ l'ens des formes \deleted{qui} sur $M$ qui sont compatibles avec $\varphi_{\underline m}|M$. Prouvons que c'est indépendant de $\underline m,\alpha$.

En effet, si on a aussi $\alpha':M\unsure{\subset|\hookrightarrow}M_{\underline m'}$, comme $\varphi_{\underline m}$ et $\varphi_{\underline m'}$ sont compatibles par hyp, il en est de $=$ de $\varphi_{\underline m}|M$ et $\varphi_{\underline m'}|M$. On voit que ces $P(M)$ forment une prépolarisation maximale, \deleted{\unsure{et}} à laquelle appartiennent \inserted{évidemment} les $\varphi_i$. Si $\varphi\in P(M),\varphi'\in P(M')$, \deleted{\unsure{qui}} on a $\varphi\sim\varphi_{\underline m}|M$, $\varphi'\sim\varphi_{\underline m'}|M'$, donc $\varphi\otimes\varphi'\sim\varphi_{\underline m}\otimes\varphi_{\underline m'}|M\otimes\unsure{M'|M)}=\varphi_{(\underline m+\underline m')}|M\otimes M'$, \deleted{\unsure{donc}}

% archive sheet number: 10
